% ---------
% --------
%!TEX TS-program = pdflatex
%!TEX encoding = UTF-8 Unicode
%!TEX spellcheck = English (Aspell)

\documentclass[11pt]{article}
\usepackage{jeffe,handout,graphicx,varwidth}
\usepackage{fourier-orns}
\usepackage{mathtools}
\usepackage[normalem]{ulem} % either use this (simple) or
\usepackage{tikz}


\def\Sym#1{\texttt{\color{BrickRed}#1}}
\allowdisplaybreaks

\begin{document}

\begin{center}
\Large\textbf{CSCI 432/532, Spring 2026}%
\\
\LARGE\textbf{Homework 9}%
\\[0.5ex]
\large Due Tuesday, March 31 at 9:00am
\end{center}

\bigskip
\hrule
\bigskip

\subsection*{Submission Requirements}
\begin{itemize}
    \item Type or clearly hand-write your solutions into a PDF format so that they are legible and professional. Submit your PDF on Canvas or Gradescope. 
	\item Put each \textbf{sub}problem on a separate page and select the corresponding problem via the Gradescope interface.
    \item Do not submit your first draft. Type or clearly re-write your solutions for your final submission. If your submission is not legible, we will ask you to resubmit.
\end{itemize}

\subsection*{Collaboration and Resources}

The \emph{best} resources for completing the homework are (in no particular order):
\begin{itemize}
	\item Your own notes from lectures and problem sessions, and/or lecture recordings.
	\item The lecture notes linked from the course website for the lecture day.
	\item The questions channel on Discord.
	\item Office hours.
	\item Discussions with classmates.
\end{itemize}
You may also access almost any resource in preparing your solution to the
homework. The exception is that you may not seek answers to the problems on the
homework directly, such as by pasting them into a GenAI tool, searching for
solutions on a search engine, looking for them on Chegg or similar websites, or
asking another person for the solution.

Additionally, you \EMPH{must}
\begin{itemize}
    \item Write your solutions in your own words---this means that you should never be \emph{copying}
    anything of substance from any source, and
    \item credit every resource you use, including GenAI tools, by providing links or full chat transcripts.
    If you use the provided LaTeX template, you can use the \texttt{Sources}
    environment for this. Ask if you need help!
\end{itemize}


\begin{boxquote}{0.9}
Remember, if you choose to use GenAI tools, not only must you provide a full
transcript of your conversation (or a link to one), you must also use the
following prompt (or some variation of it), and provide a full transcript of
the conversation.
\end{boxquote}  

Prompt:
\begin{boxquote}{0.9}
You are acting as a teaching assistant for an advanced undergraduate/graduate
algorithms and computer science theory course.

Your role is not to provide full solutions or final answers.

Instead, you should act like we are having a conversation. Ask me a single
question and let me respond. Avoid asking me many questions at once or giving
me a lot of information at one time. Guide me using hints, leading questions,
partial insights, and sanity checks.  Help me debug my own proofs, algorithms,
or reductions. Do not give a complete solution or full proofs. Be patient with
me. Don't give me details so that I can move on to the next part of a problem.
Make sure that I understood before moving on.

The course uses Jeff Erickson's Models of Computation lecture notes, so you should
guide me to answer problems using the definitions, notation, and style from those notes.

Assume I am a serious student trying to learn, not to shortcut the assignment.
I am attaching the assignment, so you should refuse to provide a complete solution
to the problems listed. Of course, you can walk me through the "Solved
Problems" listed at the end if I ask.
\end{boxquote}  

\newpage
%----------------------------------------------------------------------

\headers{CSCI 432/532}{Homework 9 (due March 31)}{Spring 2026}

Rubric for algorithms below:

\begin{Rubric}{Standard algorithm rubric}~
	\begin{itemize}
\item A complete description always includes an analysis of the algorithm's
running time, even if the problem statement does not include the word “analyze”.
(On the other hand, unless explicitly stated otherwise,you do not need to analyze
the algorithm's space usage.) 
\item Watch for hand-waving, pronouns without clear antecedents (especially
“this”), overconfidence red flags (like “simply” or “just”), meaningless filler
words (like “go through the array and”), and the
Deadly Sin “repeat this process”.
\item Do not regurgitate standard algorithms we've already seen, or algorithms
you have seen in prerequesite courses. If you use a standard algorithm as a
subroutine, you can just call that algorithm by name, but be explicit about
what it takes in and what it outputs, and give the runtime in the analysis.
\item Every algorithm must be clearly specified in English, unless the
specification is already given precisely in the problem statement. A description
of the algorithm is not enough; we want a description of what your algorithm
computes, not just an explanation of how your algorithm works.  In particular,
if the algorithm solves a more general problem than requested, you must
explicitly specify the more general problem. Similarly, if the algorithm assumes
any conditions that are not explicit in the given problem statement, you must
state those assumptions explicitly.
\item The meaning of every variable must be either clear from context (like $n$
for input size, or $i$ and $j$ for loop indices) or specified explicitly.
	\end{itemize}
\end{Rubric}

%----------------------------------------------------------------------
\begin{enumerate}


\item
Watch the following video on convolutions (and how to compute them using FFT,
though we learned that in class):
\url{https://www.youtube.com/watch?v=KuXjwB4LzSA}.

\Hint{Use this video to understand what convolutions are and how they can be
used. Then, in part (b), how can use convolutions to get
information you need? (And what is the runtime of computing a convolution using FFT?)}

\item Let $A$ and $B$ be two sets of $n$ integers.  

\begin{enumerate}
	\item Give a $O(n^2 \log n)$-time algorithm to count the distinct
	number of elements in the set $A+B = \{a+b : a \in A, b \in B\}$.
	\item Give a $O(M \log M)$-time algorithm to count the distinct number of
	elements in the set $A+B$, where $M$ is the largest possible value of $a+b$
	for $a \in A$ and $b \in B$.
\end{enumerate}
%----------------------------------------------------------------------



\end{enumerate}



\end{document}
